% =============================================================================
% FILA-Net — draft content for the organizers' Overleaf template
% (https://www.overleaf.com/read/vjztpvrnbrdh, main.tex, ACM sigconf).
%
% How to use: copy each block below into the matching place of the
% template's main.tex, replacing ONLY the gray \guide{...} text (the black
% template text must stay unchanged). Add the entries of
% report/references_extra.bib to the template's main.bib, upload
% report/figures/pipeline.pdf and the figures written by the notebook
% (artifacts/figures/*.png, artifacts/oof_export_tuned/oof_rubric_panels.png).
%
% Every \fillme{...} is a number that must come from a saved artifact of the
% final run (file given in the placeholder). Do not submit with any \fillme
% left: define it during writing as
%   \newcommand{\fillme}[1]{\textcolor{red}{[#1]}}
% and delete the definition at the end, so LaTeX fails if one is forgotten.
% Sentences about components that the ablations reject must be removed.
% =============================================================================

% ------------------------------------------------------------------ title
\title{FILA-Net: Topology-Aware Instance Reconstruction of Solar Filaments}

% ------------------------------------------------------------------ abstract (after the black sentence)
We predict, for every pixel of a $2048\times2048$ H-$\alpha$ image, a set of
dense fields---filament region, centerline, boundary, normalized distance and
axial orientation---with a ConvNeXt-T/FPN network, and reconstruct each
filament as a separate instance with centerline seeds, an
orientation-consistent endpoint graph repair and a marker-controlled
watershed. All model and post-processing choices are made on out-of-fold
predictions scored with the official Panoptic Quality (PQ), using folds that
keep each observing day together and a nested calibration. Our final
ensemble reaches an out-of-fold PQ of \fillme{oof\_metrics\_tuned.json: pq}
(95\% CI \fillme{ci95}), compared with \fillme{data\_audit.json:
inter\_annotator PQ} between human annotators.

% ------------------------------------------------------------------ introduction (after the black text)
The official score is the micro-averaged PQ over all annotator label sets,
with a match at IoU $>0.5$; a predicted filament therefore counts as a false
positive for every annotator who did not draw it. Our design follows from
this: (i) predict instances that are complete and unfragmented, because a
filament split in two pieces with IoU $\le0.5$ each yields two false positives
and one false negative; (ii) report only instances whose expected
contribution is positive. Section~\ref{sec:methodology} describes the
pipeline (Fig.~\ref{fig:pipeline}) and Section~\ref{sec:evaluation} its
out-of-fold evaluation. Code, notebook and weights:
\url{https://github.com/BFtechlab-flow/sole}.

% ------------------------------------------------------------------ methodology
\begin{figure*}[t]
  \centering
  \includegraphics[width=\textwidth]{pipeline.pdf}
  \caption{FILA-Net pipeline. Top: inference on one image. Bottom: training,
  validation and calibration, all scored with the official PQ.}
  \label{fig:pipeline}
\end{figure*}

\paragraph{Data and splits.}
The training set contains \fillme{data\_audit.json: unique\_observations}
observations with \fillme{annotation\_records} annotator label sets and
\fillme{filament\_instances} filaments. Every label set is a training sample
(observations labeled by several annotators are seen once per annotator).
We use 5 folds grouped by UT day: all observations of a day---the same Sun
seen by several GONG sites minutes apart---and all their label sets belong to
one fold, so the out-of-fold (OOF) score is not inflated by near-duplicates.
Only the annotations provided by the competition are used.

\paragraph{Inputs.}
Each grayscale image is turned into four channels: intensity normalized
between its 1st and 99th percentiles, CLAHE~\cite{zuiderveld1994clahe}, a
ridge response (difference of Gaussians, $\sigma=1.2$ and $5$) that highlights
thin dark threads, and the normalized radial position, which lets the network
account for limb darkening.

\paragraph{Network.}
A ConvNeXt-T encoder~\cite{liu2022convnext} initialized on ImageNet (via
timm~\cite{rw2019timm}) feeds a feature pyramid~\cite{lin2017fpn} fused at
stride~4, refined by a multi-dilation context block, a high-frequency branch
and a limb-aware gate (\fillme{kept/removed after ablation}), and upsampled to
full resolution. $1\times1$ heads predict the region, centerline (dilated
skeleton), boundary and per-filament normalized distance, the axial
orientation $(\cos2\theta,\sin2\theta)$ of the skeleton, and auxiliary
width, endpoint/junction and aleatoric-uncertainty maps
\cite{kendall2017uncertainties}. All targets are derived from the
segmentation polygons only.

\paragraph{Training.}
The loss combines BCE and Dice on region and centerline, BCE on the boundary,
a smooth-$L_1$ distance term, clDice~\cite{shit2021cldice} for topological
continuity, a cosine loss on the axial orientation and small weights on the
auxiliary heads (\fillme{final weights: configs/filanet\_tuned.yaml}). We
train on random $1024\times1024$ crops (70\% centered on a filament) with
flips, AdamW (lr $2\cdot10^{-4}$, cosine schedule, warm-up), mixed precision
and gradient clipping, for \fillme{epochs} epochs. Every two epochs the
held-out fold is predicted in full and scored with the official PQ; the best
epoch is kept.

\paragraph{Inference.}
Images are processed in $1024$-pixel tiles with $256$-pixel overlap blended
by a Hann window, with flip test-time augmentation (the orientation field is
transformed consistently) and averaging over the five fold models.

\paragraph{Instance reconstruction.}
The region probability is thresholded and small components are removed; the
thresholded centerline is skeletonized to obtain seeds. Skeleton endpoints of
different fragments are joined when the gap is short, aligned with the
predicted orientation on both sides, supported by the region probability,
free of boundary evidence and compatible in width (greedy acceptance in order
of cost, without cycles). A marker-controlled watershed~\cite{beucher1993watershed}
on the distance map (penalized at predicted boundaries) then assigns every
foreground pixel to one filament, giving disjoint instances; holes are filled
and instances optionally dilated. Finally each instance gets a confidence
(its mean region probability): since adding an instance with match
probability $p$ and expected IoU $u$ increases the micro PQ only if
$p\,u > \mathrm{PQ}/2$, instances below a calibrated confidence are dropped.
All reconstruction parameters are calibrated on OOF predictions
(Section~\ref{sec:evaluation}). Masks are written as COCO RLE~\cite{lin2014coco}.

\begin{figure}[t]
  \centering
  \includegraphics[width=\linewidth]{test_predictions.png}
  \caption{Predicted filaments (one color per instance) on test images.
  \fillme{choose 2--3 representative OOF examples with ground truth instead,
  including one failure case}.}
  \label{fig:examples}
\end{figure}

% ------------------------------------------------------------------ evaluation
\paragraph{Strategy.}
Our evaluation code reproduces the organizers' notebook exactly (identical
PQ on random multi-annotator scenes\fillme{ and on our OOF export: run
tools/verify\_against\_official\_notebook.py on artifacts/oof\_export}). We report OOF
scores: every training image is predicted by the model of the fold that did
not see its day. Uncertainty comes from 2{,}000 bootstrap resamples of whole
observations. Post-processing is calibrated by coordinate ascent with a
nested scheme: for each fold, parameters are chosen on the other four and
scored on it, which gives an unbiased estimate of the calibrated pipeline.

\begin{table}[t]
  \caption{Out-of-fold results (official micro PQ, 95\% bootstrap CI).}
  \label{tab:main}
  \small
  \begin{tabular}{lcccc}
    \toprule
    & PQ & SQ & RQ & TP / FP / FN \\
    \midrule
    Human (annotator vs.\ annotator) & \fillme{} & \fillme{} & \fillme{} & \fillme{} \\
    FILA-Net, default post-proc. & \fillme{} & \fillme{} & \fillme{} & \fillme{} \\
    FILA-Net, nested calibration & \fillme{} & -- & -- & -- \\
    FILA-Net, final (tuned) & \fillme{} & \fillme{} & \fillme{} & \fillme{} \\
    \bottomrule
  \end{tabular}
\end{table}

\begin{table}[t]
  \caption{Ablation (paired bootstrap of the OOF PQ difference with respect
  to the previous row; \fillme{all folds / fold 0 screening}).}
  \label{tab:ablation}
  \small
  \begin{tabular}{lcc}
    \toprule
    Variant & PQ & $\Delta$PQ [95\% CI] \\
    \midrule
    Baseline (BCE+Dice, no repair, no TTA) & \fillme{} & -- \\
    + clDice topology loss & \fillme{} & \fillme{} \\
    + orientation head & \fillme{} & \fillme{} \\
    + endpoint graph repair & \fillme{} & \fillme{} \\
    + flip TTA & \fillme{} & \fillme{} \\
    + calibrated post-processing & \fillme{} & \fillme{} \\
    \bottomrule
  \end{tabular}
\end{table}

\begin{figure}[t]
  \centering
  \includegraphics[width=\linewidth]{oof_rubric_panels.png}
  \caption{OOF distributions of IoU and Dice over overlapping
  ground-truth/prediction pairs, and one-to-many (fragmentation) /
  many-to-one (merging) relations.}
  \label{fig:rubric}
\end{figure}

\paragraph{Results.}
Table~\ref{tab:main} summarizes the OOF scores and Table~\ref{tab:ablation}
the contribution of each component. Fig.~\ref{fig:rubric} shows the
rubric distributions: \fillme{fragmentation\_rate, merge\_rate, miss\_rate,
spurious\_rate from oof\_metrics\_tuned.json}. The public leaderboard score
is \fillme{score} (used only as a sanity check). Test inference with the
5-model ensemble and TTA runs at \fillme{submission.verification.json:
images\_per\_second} images/s on \fillme{GPU}.

\paragraph{Strengths and weaknesses.}
\fillme{Write from the measured numbers. Typical items to check:} the dense
orientation and centerline fields make long curved filaments continuous;
remaining errors are \fillme{misses of faint, small filaments / splits at
gaps wider than the repair distance / merges of nearby filaments}, and part
of the error is annotator disagreement (the human baseline above).
Improvements: global (instead of greedy) endpoint matching, boundary
refinement at full resolution, and self-supervised pre-training on unlabeled
GONG frames.
